\documentclass[11pt,a4paper]{article}
\usepackage{../pelm}
\begin{document}
\answerpaper{Quiz 2}{Weeks 3--4: How LLMs Work, and Principles of Effective Prompting}

\begin{enumerate}[leftmargin=*, itemsep=0.5em, label=\textbf{\arabic*.}]
\key{B}{Tokens are not words. Common words are often one token; rarer words fragment; spaces are
  usually attached to the following word.}
\key{B}{A fixed vocabulary of whole words fails on any unseen word. Sub-word pieces mean the system
  is never stuck --- new names, technical terms and typos can all be assembled.}
\key{B}{Agglutinative stacking produces long words that appeared rarely in training, so they
  fragment. Consequences: the context fills faster, metered usage costs more, and coverage is
  thinner.}
\key{B}{Letter counting asks about something the system never perceived. This is a predictable
  consequence of tokenisation, not a general weakness at arithmetic.}
\key{B}{A probability distribution over the whole vocabulary, from which one token is selected ---
  then the process repeats. Nothing is retrieved and no rules are applied.}
\key{B}{Always selecting the most probable token produces flat, repetitive text, so randomness is
  introduced deliberately. Often exposed as ``temperature''.}
\key{B}{Consistency measures how strongly a pattern is represented, not whether it is true. A
  widely repeated error is reproduced very consistently.}
\key{C}{The system optimises plausibility. Plausible usually correlates with true, but nothing
  checks, so where they diverge nothing notices. This is why it cannot simply be patched out.}
\key{C}{Instructions change tone and shape, never available information. Asking for authority
  produces authority --- which makes a wrong answer more convincing, not less.}
\key{C}{Wrong shape means Format. Wrong content means Context or Task. Correct but lifeless means
  Examples. Format is the cheapest lever and the most neglected.}
\end{enumerate}
\end{document}
